\section{The source change that enables the additional fold}
\label{sec:source-rotation}

Let $M_{237}$ denote the composed candidate at ProveIt commit
\path{1b27f052d770f4c7d435e82e779bf21f5a1f66ed}, with inherited labels,
initial state $95$, dispatch state $96$, and halt label $254$.
Its source is the literal file \texttt{baseline\_sites.json} distributed here.
This section establishes the source part of the construction, independently of
the state-fold argument and of the binary runtime checker.

The relevant modification is one reassociation, at the source prefix
\[
 p=\texttt{1110100100110110101010011100001}.
\]
At that prefix the old source has the form $P;(Q;R)$.  Replace it by
$(P;Q);R$.  The three chunks contain respectively $245$, $160$, and $132$
instruction occurrences.  Each is syntactically closed in the precise sense
below.  There is exactly one occurrence of the modified source root.

\subsection{The elementary source identity}

A source is a finite, full binary tree with instruction leaves.  A leaf has a
binary address $a$.  Increment and positive decrement advance its binary
instruction address; a zero decrement removes the suffix through one zero and
then advances; a break of depth $k$ removes the suffix through $k$ zeros and then
advances; and a continue of depth $k$ removes the suffix through $k$ ones.
A halt leaf halts.  Subsequent lookup pads the address with zeros and stops at
its encountered instruction leaf.

Call a chunk \emph{syntactically closed} when, for every relative leaf path $a$:
\begin{itemize}
\item a zero decrement has at least one zero in $a$;
\item a break of depth $k$ has at least $k$ zeros in $a$;
\item a continue of depth $k$ has at least $k$ ones in $a$;
\item no halt leaf occurs.
\end{itemize}
These are finite syntactic checks, rather than assumptions about the numerical
values of any registers.  The supplied reconstruction script independently
checks all three chunks against these conditions.

For closed chunks, associativity
\[
 P;(Q;R)\;\simeq\;(P;Q);R
\]
holds for the structured register program.  The corresponding address changes
inside the chosen occurrence are exactly
\[
 \begin{array}{c|c|c}
 \text{chunk}&\text{old address}&\text{new address}\\ \hline
 P& p0a&p00a\\
 Q& p10a&p01a\\
 R& p11a&p1a
 \end{array}
\]
with the relative suffix $a$ and every operation and register operand unchanged.
Closure ensures that the location selected by every internal control transfer
is preserved under these changes, after the next instruction lookup.  Ordinary
fallthrough follows the same ordered instruction occurrences.  Thus each
instruction has the same successor occurrence, including both outcomes of every
decrement, and it makes the same change to the register vector.

This associativity principle is inherited background: the earlier Source~103
construction already used closed reassociations.  The contribution here is the
particular occurrence selected and its compatibility with a different decoder
sharing.  No novelty is claimed for the associativity law itself.

\subsection{An exact finite control certificate}

The construction does not rely solely on the explanatory associativity argument.
The unmodified Source~103 checker \texttt{check\_rotation\_certificate.py} was
replayed with a new, one-stage rotation certificate and the exact pinned input
bytes.  It constructs the full source tree and records each leaf in left-to-right
order.  For each instruction, it computes every applicable successor by literal
binary increment, literal suffix stripping, and zero-padded tree lookup.  It
checks that no nonzero address suffix is discarded after a leaf is reached.
Both decrement guards are included.

The old and new ordered control graphs are exactly equal.  They contain $6{,}455$
instruction occurrences, $8{,}385$ nonhalting transition descriptors, and one halt
descriptor.  Equivalently, the check covers $8{,}386$ instruction/guard or halt
cases.  Their common normalized graph has SHA-256
\begin{quote}\small\ttfamily
 ae7f12f37306b86e5f70469210fd2a8d68b6f9df094bf659abc6b4a56714f1fd.
\end{quote}
Exactly $537$ paths change: instruction indices $1501$ through $2037$, with indices
starting at zero.  All ordered operation labels and operands are identical.
The same entry leaf, index zero, corresponds to the empty initial address.
Consequently the relation that pairs equal instruction indices and equal natural
register vectors is a bisimulation of the normalized source executions.  In
particular it preserves halting for the inherited logical seed
\[
 (3,0,0,0,3,0,0,0,0).
\]
The claim quantifies over every natural register vector; it is not an observed
execution of the source from a small sample of inputs.

\subsection{Re-encoding the modified decoder}

The original decoder realizes $P;(Q;R)$ by dispatching state $177$ to state $165$
for $P$ and state $185$ for $Q;R$.  State $185$ dispatches to states $192$ and $212$
for $Q$ and $R$.  Reuse the same states with the following four row edits:
\[
 \begin{array}{c|c|c}
 \text{row}&\text{old action}&\text{new action}\\ \hline
 177/0&0\ R\ 165&0\ R\ 185\\
 177/1&1\ R\ 185&1\ R\ 212\\
 185/0&0\ R\ 192&0\ R\ 165\\
 185/1&1\ R\ 212&1\ R\ 192
 \end{array}
\]
\Needspace{6\baselineskip}
All four actions preserve the scanned bit and move right.  This leaves $237$
working states in the intermediate table \texttt{rotated237\_old.tm}.  Its SHA-256
is
\begin{quote}\small\ttfamily
 da0b3d1fdd74d56482b9d231119bbbaccdbcadc78e7b7df72c3a734a3a808724.
\end{quote}
The new literal source \texttt{rotated\_sites.json} has SHA-256
\begin{quote}\small\ttfamily
 1ebe9691507b2d3236951a31d6b9a7ef548a2cdf90525cd57051992f8ac2626f.
\end{quote}
There are still $224$ distinct source-DAG nodes, consisting of $201$ interior
nodes and $23$ instruction leaf kinds.  The longest zero run remains six.
The maximum path length increases from $48$ to $49$.

This is why an objective that only decreases source-DAG size or path depth misses
the successful modification.  The useful change is the meaning of an existing
routing state.  The resulting decoder can share two further control roles, as
established by the separate literal-table fold certificate.

Source graph equality alone does not prove a binary Turing machine correct:
backward decoder routes and cleanup routines also depend on tape contexts.
The construction therefore checks the complete changed table directly with the
unbounded symbolic runtime checker, and checks its blank startup against the
complete inherited frame.  The intermediate passes all $8{,}386$ quantified
cases, using $2{,}262{,}455$ symbolic proof steps; the minimum concrete duration
is $13$ steps.  The four edited rows are all unvisited in the inherited $409{,}835$-step
blank startup, so this source change preserves that startup step for step.
An exact sparse and dense replay of the $237$-state intermediate independently
reaches the same complete frame after $409{,}835$ steps.  These binary
checks are separate from, and complementary to, the source bisimulation proved
in this section.

\subsection{Reproduction and proof scope}

Run \texttt{python build\_and\_check\_source.py} in the supplied source-certificate
directory, without Python's optimization flags.  The script verifies the two
pinned input hashes, finds the stated source occurrence, checks chunk closure,
replays the exact one-stage control certificate, reconstructs the intermediate
literal table, and checks its expected hash.  It writes
\texttt{source\_receipt.json} with the graph identity, leaf census, source and
table hashes, source-seed binding, and row edits.

The checker is a transparent finite arithmetic computation supporting an
ordinary mathematical proof formalizable in ZFC.  It is not a new Lean or Rocq
endpoint, and its successful execution is not described as a kernel-checked
internal proof in PA.
